\section{Question 2: brevet à effet unitaire}

\subsection{Procédure}

Pour mémoire, nous sommes le 08/03/2025. R6(1) RUPC dispose que le délai 
pour déposer la demande d'effet unitaire
  est de 1 mois après la publication de la mention de la mention de délivrance
  du brevet EP au Bulletin Européen des Brevets;
donc le délai pour demander l'effet unitaire expire le 26/03/2025.
Donc, le 08/03/2025 est dans le délai pour demander l'effet unitaire.\\

Par ailleurs, R6(2) RUPC dispose que la demande d'effet unitaire 
doit être présentée par écrit dans la langue de la procédure,
c'est-à-dire en français, et doit contenir:
\begin{enumerate}
    \renewcommand{\labelenumi}{\alph{enumi})}
    \item le nom, l'adresse, la nationalité et l'adresse du siège du demandeur: 
    (son nom est James,) James Craig, MI6 Royaume-Uni, Anglais, société française;
    \item le numéro du brevet EP: EP01;
    \item le nom et l'adresse du mandataire: Aurélien, 1 rue du CEIPI;
    \item une traduction en anglais du fascicule du brevet européen
    car la langue de procédure est le français.\\
\end{enumerate}

Le brevet à effet unitaire ne couvre pas des pays au choix. En fait, le brevet à effet unitaire
couvre uniformément les États membres participants dans lesquels l'Accord JUB est en vigueur
à la date d'enregistrement de l'effet unitaire. À la date du 08/03/2025, l'effet unitaire couvre
18 états: AT, BE, BG, DK, EE, FI, FR, DE, IT, LV, LT, LU, MT, NL, PT, RO, SI, SE.\\

Or, les pays souhaités sont: FR, DE, NL, DK, NO, UK. Donc, FR, DE, NL et DK sont couverts par
l'effet unitaire. En revanche, il est nécessaire de valider le brevet EP pour NO et UK.\\

Pour valider le brevet EP pour NO et UK, il faut dans un délai de 3 mois par rapport à la date
de publication de la mention de la délivrance du brevet au Bulletin Européen des Brevets:
\begin{enumerate}
        \renewcommand{\labelenumi}{\alph{enumi})}
    \item fournir une traduction en anglais ou en norvégien du brevet pour NO
    par A1(2) du Protocole de Londres, 
    pas de traduction à fournir pour UK par A1(1) du Protocole de Londres;
    \item désigner un mandataire local si requis;
    \item payer une taxe de publication de la traduction si requise;
    \item payer les taxes annuelles.
\end{enumerate}

\subsection{Stratégie}

NL308' divulgue un instrument à commande manuelle pour seringue ("ce dispositif 100 est réalisé pour être
    adapté sur des seringues 120" par.[0007]) comportant:
\begin{itemize}
    \item un piston pourvu d'un siège adapté pour recevoir un axe de commande du coulissement
    du piston à l'intérieur du corps de la seringue ("des seringues qui intègre un piston 121"
    par.[0007]);
    \item un moyen d'accouplement au siège du piston ("un clip élastique 105 pour son couplage"
    par.[0007]);
    \item un axe fileté portant un filetage extérieur et en extrémité duquel le moyen
    d'accouplement est solidement fixé ("une première extrémité du piston fileté 104" par.[0007]);
    \item un écrou d'actionnement manuel pourvu d'un trou traversant portant un filetage
    intérieur ayant le même pas que le filetage de l'axe fileté et au travers duquel l'axe
    fileté passe et vient en prise ("le filtage extérieur du piston fileté 104 est également en prise
    avec le filetage intérieur d'une bague de dosage 108" par.[0009]);
    \item une cinématique telle que lorsque le moyen d'accouplement est couplé au siège du piston,
    la rotation de l'écrou d'actionnement par rapport au corps de la seringue provoque le coulissement
    axial de l'axe fileté au travers de son trou traversant et le déplacement du piston à l'intérieur
    du corps de la seringue ("une rotation de la bague 108 du dispositif 100 permet alors d'ajuster
    avec précision la course de l'ensemble formé par le piston fileté 104 avec le piston 121 de la
    seringue 120" par.[0009]).\\
\end{itemize}

NL308' ne divulgue pas une virole annulaire telle que définie à la revendication 2.\\

NL308' divulgue un procédé n'utilisant pas la virole annulaire définie à la revendication 3:
"lorsque le corps cylindrique 102 prend
appui contre l'extrémité du réservoir de la seringue 120, la rotation de la bague 108 permet de doser le volume
aspiré par la seringue 120 en remplissage" (par.[00010]). En revanche, NL308' ne divulgue
pas le procédé utilisant la virole annulaire.\\

Donc, la revendication 1 et 3 ne sont pas nouvelles par rapport à NL308'. Pas d'attaque d'activité
inventive possible avec NL308': cf. question 1.\\

La divulgation par NL308'
dans une action en nullité auprès de la JUB risque de réduire la portée de protection du brevet 
à effet unitaire sur tous les territoires de la JUB, notamment FR, DE, NL et DK.\\

\textbf{Conclusion:
\begin{enumerate}
    \item demander un opt-out;
    \item déposer R1+R2 et R3(virole annulaire) pour NL;
    \item déposer R1, R2, R3 pour FR, DE, DK, NO, UK.
\end{enumerate}}

